\documentclass[12pt]{article}
\usepackage[margin=1in]{geometry}
\usepackage{setspace}
\doublespacing

\title{The Warrant Question:\\Three Cases of Guarded Crossing}
\author{Flyxion}
\date{}

\begin{document}

\maketitle

\section*{}

A parish priest in a Basque-adjacent Spanish village discovers he no longer believes in the resurrection he is paid to promise. An inquisitor sent to the Basque mountains to confirm a wave of witchcraft discovers that the confessions he has been sent to ratify collapse the moment they are tested. One is fiction, drawn from Unamuno's \textit{San Manuel Bueno, m\'artir}; the other is history, drawn from the tribunal reports of Alonso de Salazar Fr\'ias. Nothing obliges these two men to be read together. But they are solving structurally the same problem, from opposite sides of it, and setting them side by side clarifies something that neither case shows fully on its own: continuation of a community does not depend primarily on what is true. It depends on discipline about what is permitted to become consequential.

Call this discipline admissibility. Not admissibility in the sense of what may be believed privately, which no one can police, but admissibility in the narrower and more consequential sense of what is allowed to cross from a private epistemic state into a public, operative one: into a verdict, a liturgy, a shared account of the world that other people will act inside. Salazar and Manuel are both, in effect, gatekeepers of this boundary. They differ only in which direction they are guarding it.

\section*{The Problem of Certification}

Salazar's problem is that testimony is trying to cross upward, from confession into fact, and he does not think it has earned the crossing. He is sent to the Logro\~no region already primed by his superiors to confirm what earlier inquisitors had found: that the confessions extracted from accused witches described a real conspiracy with the Devil, complete with nocturnal flights, sabbaths, and diabolical pacts.\cite{salazar1612} What he finds instead is that the confessions do not hold together. People report events that, on his own testing, cannot be replicated or corroborated. Different judges classify the same confession as good, bad, or indifferent with no shared method, so that the sorting tells you more about the judge than about the confessed act. 

Salazar's conclusion is not that witchcraft is metaphysically impossible; he explicitly refuses that stronger claim, since it is not the question in front of him. His conclusion is narrower and more disciplined: confession alone, however sincere, however consistent across many accused, does not constitute the ``external and objective evidence sufficient to convince everyone who hears it'' that a capital judgment requires. He is refusing to let an unverified private state, however emotionally total for the person reporting it, become load-bearing for other people's lives and deaths.

This is a harder position than either belief or disbelief in witches, because it does not resolve the underlying metaphysical question at all. It simply insists that the question of whether something happened and the question of whether a report of it is fit to govern consequences are two different questions, and that the second one can be answered even when the first cannot. Salazar's contemporaries, invested in a confirmed panic, found this maddening. It offered them no metaphysical victory in either direction. It only told them that the bridge they wanted to walk across, from a frightened child's story to a sentence of execution, was not load-bearing, whatever was or was not on the other side of it.

\section*{The Problem of Withholding}

Manuel's problem runs the other way.\cite{unamuno1930} His private discovery is not a piece of testimony arriving from outside that he must evaluate before certifying it. It is his own conclusion, arrived at from within, that the immortality he preaches every week is not something he can any longer inhabit as a live possibility. The crossing he is guarding against is downward: from his private epistemic collapse into the villagers' functioning world. For them, the Christian account of death remains reachable. Grief can still be placed inside resurrection; suffering can still be understood as temporary; the rites Manuel performs are not decorations on top of belief but the actual machinery by which death, marriage, and reconciliation stay legible to the people living through them. 

If Manuel discloses what he has concluded, he does not simply add a new fact to the villagers' stock of information. He risks making a whole range of previously reachable psychological states unreachable, the way Salazar's superiors wanted a confession to make a verdict reachable regardless of whether the underlying event had occurred.

Both men, in other words, are refusing a promotion. Salazar refuses to let confession be promoted into proof. Manuel refuses to let his own doubt be promoted into doctrine for anyone but himself. In each case the refusal is not a claim about ultimate truth. Salazar never says witchcraft is impossible; Manuel never says God does not exist. Each is making a narrower claim about what a given piece of private material is fit to carry once it leaves the person who holds it.

The asymmetry between them is worth sitting with, because it shows that admissibility is not a single fixed rule but a discipline that has to be applied freshly to each direction of crossing. Salazar's caution protects the accused from the state's willingness to treat unreliable testimony as sufficient warrant for violence. Withholding certification here is protective almost by definition: it stops something from happening to people who have not chosen it. Manuel's caution is more ambiguous, because withholding disclosure also stops something from happening to people who have not chosen it, but the something in this case is knowledge, not punishment. The villagers are denied information that is, by any ordinary standard, relevant to decisions about how to live and how to die. Manuel's silence protects their consolation at the cost of their autonomy, in a way that has no clean analogue on Salazar's side of the boundary. 

Salazar is not deciding for the accused whether they may know something about themselves; he is deciding, on behalf of the tribunal, whether a report is fit to justify a verdict. Manuel is deciding, on behalf of an entire village, which of their own psychological states remain open to them. The first is a judgment about evidentiary weight. The second is a judgment about other people's inner lives, made without their knowledge that the judgment is being made at all.

This is where the paternalism in Manuel's position becomes visible in a way it does not for Salazar. Manuel has no verified basis for his conviction that the villagers cannot survive learning what he has concluded; that conviction is itself an unfalsified prediction, exactly the kind of unearned crossing that Salazar would refuse to certify if someone brought it to him as testimony. Manuel would not accept ``the priest is certain the village would collapse'' as sufficient external evidence for anything, if it were presented to him as someone else's claim rather than his own settled intuition. He exempts his own prediction from the very standard he implicitly holds everyone else's testimony to. Salazar, by contrast, applies the same discipline to himself that he applies to the tribunal: he does not trust his own prior expectations about witchcraft any more than he trusts the confessions, and he changes his conclusion when the material does not hold up, in public, across nine separate reports over more than a decade.

\section*{The Third Case: Unlicensed Crossing}

There is a third case worth setting beside these two, because it tests the same boundary from an angle neither Salazar nor Manuel occupies. Both of them are people refusing to certify a crossing. The friar Salvatore, in Umberto Eco's \textit{The Name of the Rose}, is someone who crosses constantly, and is understood, without holding any of the credentials that are supposed to make crossing possible in the first place.\cite{eco1980} He communicates in something drawn from Latin, from vernacular dialects, from whatever fragments of other people's speech have passed through a long and disordered life, recombined on the spot according to what the moment requires rather than according to any single grammar shared with his listeners. 

A language, properly speaking, depends on a fixed convention: a community that has agreed a given sound stands for a given thing, so that no one may swap the sign for another on private authority. Salvatore's speech fails this condition at every point. It is stitched from the broken pieces of sentences heard somewhere else, deployed as needed, correct in no one system. And yet it works. Meaning arrives. He is understood, however each listener manages the recovery, and no one involved is being deceived about anything; they are only doing what listeners always do, reconstructing intention from a signal that has no official warrant behind it at all.

This matters for placing Salazar and Manuel correctly, because it shows that neither man's discipline is a general suspicion of unlicensed communication as such. Function without proper form is not intrinsically dangerous. Salvatore's patchwork speech harms no one, and nothing in the story requires that it be corrected or certified before anyone is allowed to act on it. What Salazar and Manuel are each guarding is narrower than that: not whether meaning may cross without a fixed underlying code, but whether a specific kind of consequence, an execution, or an entire village's framework for surviving its own deaths, may be built on a crossing that has not earned the weight it is being asked to carry. 

Salvatore's \textit{disiecta membra}, fragments of other sentences reassembled for the occasion, constitute a working communicative act even though no single system authorizes them. But a confession extracted under torture is also built from \textit{disiecta membra}: pieces of local rumor, a neighbor's story, whatever the interrogator has already suggested, stitched under duress into something that resembles testimony. Salazar's insight is that the resemblance is exactly the danger. Salvatore's listeners have nothing at stake in mistaking his improvised speech for a real language; a misunderstanding costs an errand, not a life. The tribunal's listeners have everything at stake, because their reassembled confession is being asked to authorize something that cannot be undone. The same operation, fragments recombined into something that functions as if unified, is harmless in one setting and lethal in the other, and the difference has nothing to do with the coherence of the fragments themselves. It has to do with what the reassembled thing is afterward permitted to justify.

Manuel's own preaching is closer to Salvatore's condition than it first appears. By the time the novella finds him, his liturgy is no longer underwritten by the single settled conviction that is supposed to authorize it; it survives instead as a performance built from parts, doctrine recited and administered without the governing belief that would once have made it one thing rather than an assembly of pieces. And it still produces real uptake. The villagers are comforted, reconciled, prepared for death, exactly as Salvatore's listeners are given real and usable meaning by a speech with no single grammar behind it. The difference from Salvatore is that Manuel knows the gap is there and chooses to keep speaking across it anyway, which is what turns his case from a curiosity of language into a moral problem. Salvatore has nothing to withhold, because he is not aware of failing any standard. Manuel is aware, at every Mass, of exactly what standard his own performance no longer meets.

\section*{Distinctions That Cannot Be Collapsed}

There is a further distinction all three cases make necessary: the difference between what actually happened and what a record of it can later verify. Salazar is careful to keep these apart. That a child reported flying to a sabbath does not mean the child flew, but it also does not mean the report was invented from nothing; something happened to produce the account, even if that something was fear, suggestion, or the shape of the questions being asked. The event and its interpretation are not the same object, and collapsing them in either direction, treating every confession as either literally true or entirely fabricated, is a failure of exactly the discipline Salazar is trying to practice. 

The same separation matters for Manuel. That he privately doubts resurrection does not retroactively make the reconciliations, the comfort, the actual pastoral work he performed for decades fictional. Those events have their own provenance and remain real regardless of what turns out to be true about the doctrine that motivated them. Conversely, the villagers' conviction that Manuel is a saint does not establish that he possessed the interior faith they attribute to him. Usefulness does not certify truth in either direction, and neither does truth erase usefulness once it is called into question. A representation can do real causal work in the world, real reconciliations, real changes in behavior, real deaths made bearable, without that work constituting evidence for the representation's literal content, and without doubt about the content erasing the work already done.

\section*{Holding the Line}

What Salazar, Manuel, and Salvatore share, finally, is a recognition that predates all three of their moments and outlives them: that a community's capacity to keep functioning is not the same question as whether its representations are verified or its speech is properly formed, and that the boundary between private states and shared, operative ones has to be actively guarded rather than assumed to take care of itself. Salvatore shows the boundary can be crossed constantly, by an utterance certified by no system at all, with no damage whatsoever, because nothing important is being asked to rest on the crossing. Left unguarded where something important is being asked to rest on it, the same kind of fragmentary, patched-together material becomes sufficient warrant for burning people who did nothing, or for dismantling everyone else's way of surviving their own losses. 

Salazar spent a career insisting that confession is not proof. Manuel spent a life insisting, through silence rather than argument, that private certainty is not doctrine. Neither man resolved the deeper question underneath his refusal. Salazar never settled whether witchcraft was real; Manuel never settled whether God existed. What each of them did instead was harder and, in its way, more useful than an answer: he held the line on what was allowed to cross, so that the people on the other side of it could keep living inside a world that still made sense, whether or not it was ever, in the end, true.

\newpage

\begin{thebibliography}{99}

\bibitem{salazar1612}
Salazar Fr\'ias, Alonso de. \textit{Tribunal Reports to the Suprema}. 1612--1623. Spanish Inquisition Archives.

\bibitem{unamuno1930}
Unamuno, Miguel de. \textit{San Manuel Bueno, m\'artir}. 1930.

\bibitem{eco1980}
Eco, Umberto. \textit{The Name of the Rose}. Translated by William Weaver. Harcourt Brace Jovanovich, 1983.

\end{thebibliography}

\end{document}
